% =============================================================================
% Aplicaciones Reales
% =============================================================================

\section{Aplicaciones Reales}

% Casos concretos de implementación
\subsection{Casos de Implementación}

La implementación de CNN en el entorno clínico cuenta con evidencia documentada en múltiples especialidades \parencite{Laurent2025}:

\subsubsection{Viz.ai — Detección de Accidentes Cerebrovasculares}

Con 13 algoritmos autorizados por la FDA y despliegue en más de 1,600 hospitales, Viz.ai utiliza CNN para detectar oclusiones de grandes vasos con un AUC superior a 0.90, reduciendo el tiempo de inicio del tratamiento en aproximadamente 66 minutos.

\subsubsection{Mirai — Predicción de Riesgo de Cáncer de Mama}

Desarrollado por el MIT Jameel Clinic, este modelo basado en CNN predice el riesgo de cáncer de mama a 5 años a partir de mamografías, con un C-index entre 0.69 y 0.78 validado prospectivamente en 5 países. Su uso en el Mount Auburn Hospital ha permitido personalizar los intervalos de cribado.

\subsubsection{Aidoc — Hemorragia Intracraneal}

El sistema Aidoc emplea CNN para la detección de hemorragia intracraneal aguda con sensibilidad superior al 90\%, funcionando como triaje automático que prioriza casos críticos en la lista de trabajo del radiólogo.

\subsubsection{Qure.ai — Nódulos Pulmonares}

La plataforma qCT-Lung, autorizada por la FDA, utiliza CNN para detectar y segmentar nódulos pulmonares en tomografías de tórax, facilitando el cribado de cáncer de pulmón a gran escala.

\subsubsection{Shockmatrix — Triaje en Trauma}

En un ensayo con 1,292 casos de trauma en Grenoble, Francia, el sistema basado en CNN falló en 20 pacientes que los médicos detectaron, mientras que los médicos fallaron en 21 que la IA identificó, demostrando el valor complementario de la IA como herramienta de apoyo.
